%% GRAMSCI v2 paper: merge-walk + parity 4PCF + Python interface + GPU
%% Reconstructed from the June 2026 gramsci_mergewalk draft (source was on
%% laptop; PDF transcribed) and extended with the GPU implementation,
%% DESI demonstration, and Quijote validation/cosmology sections.
%%
%% Build: pdflatex gramsci_v2; bibtex gramsci_v2; pdflatex x2
\documentclass[twocolumn]{aastex701}

\usepackage{amsmath}
\usepackage{listings}
\usepackage{xcolor}
\lstset{
  basicstyle=\ttfamily\scriptsize,
  keywordstyle=\color{blue},
  commentstyle=\color{gray},
  breaklines=true,
  frame=single,
  language=Python
}

\newcommand{\rmax}{r_{\mathrm{max}}}
\newcommand{\Mpch}{\,h^{-1}\mathrm{Mpc}}
\newcommand{\todo}[1]{\textcolor{red}{[TODO: #1]}}

\shorttitle{GRAMSCI: Graph-based N-point statistics}
\shortauthors{Sabiu et al.}

\begin{document}

\title{Graph Database Solution for Higher Order Spatial Statistics in the
  Era of Big Data: GPU Acceleration and Out-of-Core N-point Functions}

\author{Cristiano G. Sabiu}
\affiliation{Department of Physics, University of Seoul, Seoul 02504,
  Republic of Korea}
\email{csabiu@gmail.com}

\begin{abstract}
We present GRAMSCI (GRAph Made Statistics for Cosmological Information), an
algorithm and open-source Fortran code for the fast computation of the
general $N$-point spatial correlation function of any discrete point set
embedded in $\mathbb{R}^n$.  Using KD-trees and a graph database, we count
all possible $N$-tuples in binned configurations within a given length
scale.

\textit{Version 2 update.}  We describe several additions to the original
code.  First, we replace the binary-search inner loop --- which cost
$O(m\log m)$ per hub--spoke pair, where $m$ is the mean neighbour count ---
with a \textit{merge-walk}: a two-pointer sorted-set intersection that
exploits the pre-sorted adjacency lists and reduces the inner loop to
$O(m)$.  Second, we implement a parity-decomposed 4PCF that separates the
signal into even and odd channels, enabling direct tests of parity
violation in the galaxy distribution.  Third, we provide a Python interface
so the Fortran engine can be called directly from NumPy arrays.  Finally,
and principally, we present a GPU port of the full query engine (OpenACC):
the 3PCF, equilateral 3PCF, 4PCF, and parity-decomposed 4PCF kernels run on
a single consumer GPU with measured speedups of $2.6\times$ (3PCF) to
$9\times$ (4PCF) over a 64-thread CPU node, and an out-of-core tiling
scheme allows graphs far exceeding device memory --- we measure a
$9\times10^9$-edge BAO-scale 3PCF on a 24\,GB card with ${\sim}20\%$
overhead.  We validate the code against its CPU reference and on Quijote
$N$-body mocks, demonstrate BAO-scale 2- and 3-point measurements on the
DESI DR1 LRG sample, and illustrate the cosmological information content of
configuration-space $N$-point functions across the Quijote parameter
variations.  The code is available at
\url{https://github.com/csabiu/Gramsci}.
\end{abstract}

\keywords{cosmology: large-scale structure --- methods: statistical ---
  methods: numerical --- galaxies: statistics}

%% ===================================================================
\section{Introduction}
\label{sec:intro}

The spatial distribution of galaxies encodes information about the
primordial density field and its subsequent gravitational evolution.  The
two-point correlation function (2PCF) has been the workhorse of this
programme \citep[e.g.][]{anderson14}, but higher-order statistics carry
complementary information that is inaccessible to the power spectrum.  The
3-point correlation function (3PCF) probes the bispectrum, which constrains
primordial non-Gaussianity and galaxy bias
\citep{sefusatti06,gaztanaga05}.  The connected 4-point correlation
function (4PCF) is related to the primordial trispectrum and provides
additional leverage on inflationary models \citep{sefusatti07}.

Computing $N$-point functions from large catalogues is expensive.  A
brute-force approach requires $O(N^k)$ operations for a $k$-point function.
Tree-based methods reduce this to $O(Nm^{k-1})$, where $m$ is the mean
number of neighbours within the maximum scale $\rmax$, but the $m^{k-1}$
factor still grows rapidly.

Several codes are publicly available.  \textsc{TreeCorr} \citep{jarvis04}
provides optimised 2PCF and 3PCF estimators for projected and
three-dimensional catalogues.  \textsc{Corrfunc} \citep{sinha20} uses SIMD
vectorisation (AVX2/AVX-512) for maximum 2PCF throughput.  \textsc{Encore}
\citep{philcox22} expands the $N$-PCF in a spherical-harmonic basis,
replacing direct enumeration with matrix products; this achieves favourable
scaling at high neighbour density but does not directly yield a
configuration-space connected 4PCF with parity decomposition.

GRAMSCI \citep{sabiu19} uses a graph-database approach: it builds a compact
sorted adjacency structure once from a KD-tree, then queries it for any
combination of 2-, 3-, and 4-point statistics without revisiting the
original coordinates.  This paper extends the original work in four
directions (Section~\ref{sec:new}): an $O(m)$ merge-walk inner kernel
(Section~\ref{sec:mergewalk}), a parity-decomposed 4PCF
(Section~\ref{sec:parity}), a Python interface (Section~\ref{sec:python}),
and a GPU port of the complete query engine with out-of-core support for
graphs exceeding device memory (Section~\ref{sec:gpu}).

%% ===================================================================
\section{The GRAMSCI Algorithm}
\label{sec:algorithm}

We summarise the algorithm here; full details are in \citet{sabiu19}.

\subsection{Graph construction}
\label{sec:graph}

Given a catalogue of $N$ objects, GRAMSCI builds a KD-tree and finds, for
each object $i$, all neighbours within $\rmax$.  These are stored as a
sorted adjacency list:
\begin{equation}
  \mathcal{N}(i) = \{j_1, j_2, \ldots, j_{m_i}\}, \quad
  j_1 < j_2 < \cdots < j_{m_i}.
  \label{eq:adjacency}
\end{equation}
Alongside each neighbour index, the radial-bin index (an integer byte) and
optionally a direction-pixel byte (for the parity 4PCF) are stored.  The
KD-tree and raw coordinates are then discarded; all subsequent NPCF queries
operate entirely on the adjacency lists.

The total memory footprint is $O(Nm)$ rather than $O(N^2)$.  The sorted
order in Equation~(\ref{eq:adjacency}) is a key invariant exploited by the
merge-walk kernel and the GPU binary-search kernels alike.

\subsection{Triangle and tetrahedron enumeration}
\label{sec:enumeration}

For the 3PCF, GRAMSCI uses a hub-and-spoke loop: for each hub $i$ and each
neighbour $j \in \mathcal{N}(i)$, it checks whether a second neighbour
$k \in \mathcal{N}(i)$ with $k > j$ is also in $\mathcal{N}(j)$.  If so,
the triangle $(i,j,k)$ is counted in the appropriate
$(r_{ij}, r_{ik}, r_{jk})$ bin.  The 4PCF extends this to tetrahedra
$(i, j_1, j_2, j_3)$ where every pair is connected.  The original
membership test was a binary search on $\mathcal{N}(j)$, costing $O(\log
m)$ per candidate.

%% ===================================================================
\section{New Contributions}
\label{sec:new}

\subsection{Merge-walk: $O(m)$ sorted-set intersection}
\label{sec:mergewalk}

Since the hub tail $\{k \in \mathcal{N}(i) : k > j\}$ and $\mathcal{N}(j)$
are both sorted by node index (Equation~\ref{eq:adjacency}), their
intersection can be found in a single linear pass with two pointers $a$
and $b$:
\begin{itemize}
\item $a$ walks the hub tail, starting at the entry after $j$.
\item $b$ walks $\mathcal{N}(j)$ from its beginning.
\end{itemize}
At each step:
\begin{itemize}
\item $\mathcal{N}(i)[a] = \mathcal{N}(j)[b]$: triangle found; accumulate
  the $(r_{ij}, r_{ik}, r_{jk})$ bin contribution; advance both.
\item $\mathcal{N}(i)[a] < \mathcal{N}(j)[b]$: advance $a$.
\item otherwise: advance $b$.
\end{itemize}
Each list is visited at most once, giving $O(m_i + m_j)$ cost per spoke ---
an $O(\log m)$ improvement over binary search.

\subsubsection{Three-way merge for the 4PCF}
\label{sec:threeway}

For the 4PCF, once the k2 walk has located a second spoke $j_2$, the k3
loop must find a third spoke $j_3$ that appears simultaneously in the hub
tail beyond $j_2$, in $\mathcal{N}(j_1)$, and in $\mathcal{N}(j_2)$.  We
generalise the two-pointer walk to three pointers $\alpha, \beta, \gamma$:
\begin{itemize}
\item $\alpha$: hub tail beyond $j_2$.
\item $\beta$: $\mathcal{N}(j_1)$, starting at the position just after
  $j_2$ (reused from the k2 walk, so no rescanning).
\item $\gamma$: $\mathcal{N}(j_2)$, pre-advanced past $j_2$'s own index
  (entries $\le j_2$ can never match the hub tail).
\end{itemize}
Let $h_\alpha, h_\beta, h_\gamma$ be the current candidate node indices.
If all three are equal, we record the tetrahedron and advance all three.
Otherwise every pointer at the current minimum is advanced:
\begin{equation}
\begin{aligned}
  h_\alpha \le h_\beta \text{ and } h_\alpha \le h_\gamma &: \quad
    \alpha \leftarrow \alpha + 1 \\
  h_\beta \le h_\alpha \text{ and } h_\beta \le h_\gamma &: \quad
    \beta \leftarrow \beta + 1 \\
  h_\gamma \le h_\alpha \text{ and } h_\gamma \le h_\beta &: \quad
    \gamma \leftarrow \gamma + 1
\end{aligned}
\label{eq:threeway}
\end{equation}
(Multiple conditions can be true simultaneously when two minima are tied;
in that case multiple pointers are advanced in the same step.)  This
visits each list at most once, giving $O(m)$ cost per $(i, j_1, j_2)$
triple compared with two $O(\log m)$ binary searches in the original code.

\subsection{Parity-decomposed 4PCF}
\label{sec:parity}

A standard 4PCF estimator gives one number per tetrahedron configuration.
GRAMSCI v2 adds a parity decomposition that separates the signal into
parity-even ($\zeta^{+}$) and parity-odd ($\zeta^{-}$) channels.  For each
tetrahedron $(i, j_1, j_2, j_3)$, the handedness is determined by the
signed volume
\begin{equation}
  V = (\mathbf{r}_{ij_1} \times \mathbf{r}_{ij_2}) \cdot \mathbf{r}_{ij_3}.
  \label{eq:volume}
\end{equation}
The direction of each spoke is also quantised into a pixel on the unit
sphere (stored as a single byte per edge at graph-construction time with
${\sim}17\%$ memory overhead), enabling efficient grouping by parity at
query time without retaining the original coordinates.

Tetrahedra whose pixelised direction vectors give a mathematically zero
signed volume (repeated pixels, coplanar pixel triples) have no defined
chirality.  These are common --- repeated pixels alone account for a few
per cent of all tetrahedra at typical pixelisations, comparable to the
amplitude of the odd channel itself --- and assigning them a sign from the
floating-point residue of Equation~(\ref{eq:volume}) would contaminate
$\zeta^{-}$ with compiler- and hardware-dependent noise.  GRAMSCI therefore
assigns degenerate tetrahedra ($|V| < 10^{-9}$) zero weight in the odd
channel; the even channel is unaffected.

The parity-odd 4PCF $\zeta^{-}$ vanishes for any statistically
parity-symmetric field; a non-zero measurement signals parity violation.
Searches for such a signal in BOSS galaxy catalogues \citep{hou23} have
generated significant interest, and independent configuration-space
measurements provide a direct cross-check on spherical-harmonic-based
analyses.

\subsection{Python interface}
\label{sec:python}

A lightweight Python package (requiring only NumPy) wraps the compiled
Fortran binary.  Input/output is handled automatically through a temporary
directory; the user supplies NumPy arrays and receives structured result
objects:

\begin{lstlisting}[caption={Minimal Python usage.},label=lst:python]
import gramsci, numpy as np

# pos: (N,3) float64,  weights: (N,) float64
r3 = gramsci.compute_3pcf(
        pos_data, w_data, pos_rand, w_rand,
        rmin=1.0, rmax=30.0, nbins=6)
print(r3.zeta)   # shape (n_configs,)
print(r3.NNN, r3.r1, r3.r2, r3.r3)

r4 = gramsci.compute_4pcf(
        pos_data, w_data, pos_rand, w_rand,
        rmin=1.0, rmax=30.0, nbins=3,
        parity=True)
print(r4.zeta_even)  # parity-even 4PCF
print(r4.zeta_odd)   # parity-odd 4PCF
\end{lstlisting}

The package is installable with \texttt{pip install -e .} from the
\texttt{python/} subdirectory of the repository.

\subsection{GPU implementation}
\label{sec:gpu}

The dominant cost in every GRAMSCI query is the enumeration loop over
hub--neighbour tuples --- millions of independent hubs, each performing
$O(m^2)$ to $O(m^3)$ work on small sorted integer lists.  This structure
maps naturally onto a GPU, and we have ported the complete query engine
(3PCF, equilateral 3PCF, 4PCF, and parity 4PCF) to OpenACC, compiled with
\textsc{nvfortran}.  The GPU binary, \texttt{gramsci\_gpu}, accepts the
same command-line options and produces identical output formats to the CPU
code.  Graph construction remains on the CPU (OpenMP) and is typically a
small fraction of the total runtime.

\subsubsection{Data layout}
\label{sec:csr}

The CPU code stores adjacency lists as per-node allocatable arrays, a
structure that cannot be mapped to device memory.  The GPU path flattens
the graph into a compressed-sparse-row (CSR) form: a 64-bit offset array
$\mathtt{ptr}(1{:}N{+}1)$ and flat edge arrays holding the neighbour index
(32-bit) and radial-bin index (8-bit) --- 5 bytes per edge, plus a
direction-pixel byte for parity queries.  Edge counts routinely exceed
$2^{31}$ for BAO-scale work (Section~\ref{sec:performance}), so all edge
indexing is 64-bit.  During the flattening, each jagged row is freed as it
is copied (with the freed pages explicitly returned to the operating
system), holding peak host memory to approximately one live copy of the
graph.

\subsubsection{Kernel design}
\label{sec:kernels}

Three design rules, found empirically, determine GPU performance
(Appendix~\ref{app:openacc} records the failure modes of the alternatives):

\textit{One thread block per hub, vector lanes on the inner loop.}  Each
hub is assigned to one OpenACC gang (CUDA thread block); the innermost
neighbour loop is spread across the gang's vector lanes.  For a fixed
hub--spoke pair, all lanes then search or read the \emph{same} adjacency
list, so a warp's memory traffic is confined to one small, cache-resident
segment.  Letting the compiler choose the mapping instead places a
different hub on every vector lane, producing fully divergent memory
access; we measured this na\"ive mapping to be $6\times$ slower than the
64-thread CPU code, versus $2.6\times$ faster for the explicit mapping.

\textit{Per-hub scratch via gang slots.}  The 4PCF kernels precompute, per
hub, a local connectivity matrix $\mathtt{lmat}(k_2,k_1)$ holding the bin
index of the edge between neighbours $k_1<k_2$, reducing the innermost
tetrahedron loop to $O(1)$ array lookups: $C(m,2)$ searches replace the
original $3\,C(m,3)$.  This matrix (up to several hundred kilobytes for
dense graphs) cannot be a private array --- OpenACC privatises arrays per
vector lane --- so each gang owns a slice of a global scratch array indexed
by an explicit gang-slot loop.  The scratch extent is sized at runtime from
the longest adjacency row of the actual graph.

\textit{Atomic accumulation into slot-strided partials.}  Vector lanes
within a gang race on shared accumulators, and compiler-generated array
reductions across nested gang/vector loops proved unreliable or slow.
Counts are instead accumulated with hardware floating-point atomics into
per-slot partial arrays (one slot per gang, or hash-strided over $\sim$4k
slots), which are summed on the host.  Atomic contention is negligible
because concurrently resident gangs own distinct slots.

\subsubsection{Out-of-core operation}
\label{sec:chunking}

At BAO scales a survey-sized graph exceeds any single GPU's memory: the
DESI DR1 LRG North+South sample at $\rmax = 150\Mpch$ produces
$9.0\times10^9$ edges (45\,GB of CSR).  The query kernels therefore
support a chunked mode, selected automatically at runtime when the graph
exceeds the free device memory.

The kernel needs two kinds of row access: the hub's own row, and the row
of any neighbour being searched.  The edge arrays are split into $W$
row-aligned windows, and the query runs as a tile loop over (hub-window,
search-window) combinations with only the active windows resident on the
device.  A triangle at hub $i$ with spoke pair $(j_1, j_2)$ is processed
exactly once --- in the tile whose hub window holds $i$'s row and whose
search window holds $j_1$'s row --- and tiles skip foreign spokes with a
single integer comparison.  For the 4PCF kernels, whose connectivity
matrix draws on two different neighbour rows, tiles run over \emph{pairs}
of search windows $(c_1 \le c_2)$; the sorted adjacency invariant
guarantees each tetrahedron lands in exactly one tile.  Window sizes are
derived at runtime from free device memory and can be overridden through
an environment variable, which also provides a convenient way to exercise
the chunked code paths on small test data.

The measured cost of chunking is modest: the $9.0\times10^9$-edge DESI
measurement above ran in 28 minutes on a 24\,GB RTX 3090~Ti, roughly 20\%
slower than the extrapolated cost of a hypothetical all-resident run, the
difference being PCIe re-transfers of the windows.

%% ===================================================================
\section{Performance Benchmarks}
\label{sec:performance}

\subsection{Test catalogue}

All CPU benchmarks used a Gaussian mock with $N_g = 207\,246$ galaxies and
$N_r = 500\,000$ randoms in a periodic box.  Timings are wall-clock
seconds on a single node with 4 OpenMP threads, divided by thread count to
give effective single-core seconds for fair comparison.  Radial bins are
equally spaced from $r_\mathrm{min} = 1\Mpch$ to $\rmax$.

\subsection{Merge-walk speedup}

Table~\ref{tab:mergewalk} summarises the CPU results.  The speedup grows
with $\rmax$ because the mean neighbour count $m \propto \rmax^3$ and the
savings per inner-loop step scale as $\log_2 m$.  At $\rmax = 55\Mpch$
($m \approx 700$) the 3PCF achieves $2\times$; at $\rmax = 40\Mpch$ the
4PCF achieves $3.8\times$ because two independent binary searches were
replaced by a single three-pointer walk.

\begin{deluxetable}{lrrrr}
\tablecaption{Wall-clock speedup of the merge-walk over binary search.
\label{tab:mergewalk}}
\tablehead{
  \colhead{Kernel} & \colhead{$\rmax$} & \colhead{Bsearch} &
  \colhead{Merge-walk} & \colhead{Speedup} \\
  \colhead{} & \colhead{[$h^{-1}$Mpc]} & \colhead{[s]} &
  \colhead{[s]} & \colhead{}
}
\startdata
3PCF          & 30 &   6.6 &   4.5 & 1.46$\times$ \\
3PCF          & 40 &    34 &    21 & 1.65$\times$ \\
3PCF          & 55 &   223 &   109 & 2.04$\times$ \\
4PCF          & 30 &  61.6 &  18.8 & 3.27$\times$ \\
4PCF          & 40 & 666.7 & 174.1 & 3.83$\times$ \\
4PCF (parity) & 30 &  69.2 &  24.8 & 2.79$\times$ \\
\enddata
\tablecomments{Single-node, 4 OpenMP threads; times divided by thread
  count.  The 4PCF gain is larger than the 3PCF gain because the original
  code had two binary searches in the innermost loop.}
\end{deluxetable}

\subsection{GPU speedup}
\label{sec:gpubench}

Table~\ref{tab:gpu} compares the GPU query kernels against the full
64-thread CPU node (not single-core equivalents) on the same test
catalogue and on the DESI DR1 LRG sample (Section~\ref{sec:desi}).  The
GPU advantage grows with the order of the statistic: the 3PCF reaches
$2.6\times$ over 64 CPU threads, while the 4PCF --- whose $O(m^3)$ inner
loop amortises the per-hub setup --- reaches $9.0\times$.  In all cases
the binned counts agree exactly with the CPU reference.

\begin{deluxetable}{llrrr}
\tablecaption{GPU query times vs.\ a 64-thread CPU node (RTX 3090~Ti
  vs.\ 64 OpenMP threads).\label{tab:gpu}}
\tablehead{
  \colhead{Kernel} & \colhead{Catalogue} & \colhead{CPU$\times$64} &
  \colhead{GPU} & \colhead{Speedup} \\
  \colhead{} & \colhead{} & \colhead{[s]} & \colhead{[s]} & \colhead{}
}
\startdata
3PCF, $\rmax{=}80$    & test (0.7M)    &   9.4 &   3.6 & 2.6$\times$ \\
4PCF, $\rmax{=}50$    & test (0.7M)    &  15.7 &   2.9 & 5.5$\times$ \\
4PCF, $\rmax{=}65$    & test (0.7M)    & 146.8 &  16.4 & 9.0$\times$ \\
3PCF, $\rmax{=}120$   & DESI LRG (3.0M) & \nodata & 380 & \nodata \\
3PCF, $\rmax{=}150$\tablenotemark{a} & DESI LRG (3.0M) & \nodata & 1674 & \nodata \\
\enddata
\tablenotetext{a}{Out-of-core: $9.0\times10^9$ edges (45\,GB CSR) on a
  24\,GB device, $5\times5$ window tiles.}
\tablecomments{Query time only; graph construction (CPU, OpenMP) adds
  $\lesssim 10\%$ for the GPU runs. \todo{re-run CPU reference timings for
  the DESI rows on the 64-thread node for completeness}}
\end{deluxetable}

\subsection{Correctness verification}
\label{sec:correctness}

Both implementations are run on the same catalogue and compared bin by
bin.  All four GPU kernels (3PCF, equilateral 3PCF, 4PCF, parity 4PCF)
were validated against the CPU reference in both single-pass and chunked
modes --- eight combinations --- with exact agreement of all binned counts.
The regression test suite (included in the repository) checks
isolated-pair 2PCF, regular-tetrahedra connected 4PCF, and
chiral-tetrahedra parity 4PCF against known analytic values, and
\texttt{tests/benchmark\_3pcf.py} automates the CPU--GPU comparison.

%% ===================================================================
\section{Demonstration on DESI DR1 LRG}
\label{sec:desi}

\todo{tighten this section once EZmock error bars are final}

As a realistic end-to-end demonstration we measure the 2PCF and 3PCF of
the DESI DR1 LRG sample ($1.48$M galaxies NGC, $0.66$M SGC; weights
$w \times w_{\rm FKP}$; DESI fiducial cosmology) with randoms subsampled
1:1.  Galactic caps are measured independently and combined at the level
of raw counts (weighting each cap's normalised counts by the cube --- or
square, for pair counts --- of its summed data weights).

The combined $r^2\xi(r)$ shows the acoustic peak at
$r \simeq 97.5\Mpch$ with the expected dip--peak structure, recovered
independently in both caps.  The 3PCF, measured to $\rmax = 150\Mpch$ in
5$\Mpch$ bins (the out-of-core measurement of
Table~\ref{tab:gpu}), shows the corresponding BAO feature: in the
isoceles slice $\zeta(52, 78, r_3)$ the acoustic bump appears at
$r_3 \simeq 107$--$118\Mpch$, an order of magnitude above the
surrounding baseline, with both Galactic caps showing the rise
independently.  Per-configuration uncertainties are estimated from 25 (of
1000) EZmock realisations processed through the identical pipeline.

\todo{Figure: 2PCF NGC/SGC/combined $r^2\xi$.  Figure: 3PCF slice with
EZmock error bars.  Quote the measured runtimes per mock.}

%% ===================================================================
\section{Validation and Cosmology Dependence with Quijote}
\label{sec:quijote}

\todo{This section's measurements are planned; pipeline is in place.}

\subsection{Setup}

We use Friends-of-Friends halo catalogues from the Quijote simulation
suite \citep{villaescusa20}: $512^3$ particles in $(1\,h^{-1}{\rm
Gpc})^3$ boxes, with the fiducial cosmology and the single-parameter
variations ($\Omega_{\rm m}^{\pm}$, $\sigma_8^{\pm}$, $h^{\pm}$,
$n_s^{\pm}$).  A fixed number-density cut ($\bar n = \todo{value}$) is
applied to each catalogue, and randoms are generated uniformly within the
periodic box at 1:1.

\subsection{Validation against the fiducial cosmology}

\todo{3PCF of fiducial Quijote haloes vs.\ tree-level perturbation theory
prediction at large scales; GPU vs.\ CPU agreement (exact) and timing at
box scale.  N realisations for the error bars.}

\subsection{Response of the 3PCF and 4PCF to cosmological parameters}

\todo{For each varied parameter: measure 3PCF (20 bins to 150 Mpc/h) and
4PCF (coarser bins) on $\mathcal{O}(50)$ realisations per cosmology;
plot the fractional response $(\zeta_{\theta^+} -
\zeta_{\theta^-})/\zeta_{\rm fid}$ per configuration with fiducial-mock
error bars; discuss which configurations carry the strongest parameter
sensitivity, and contrast with the 2PCF response.}

%% ===================================================================
\section{Scientific Applications}
\label{sec:applications}

\paragraph{Primordial non-Gaussianity}
The connected 4PCF $\zeta_{\rm conn}$ is a direct configuration-space
probe of the primordial trispectrum.  The $9\times$ GPU speedup at
BAO-relevant scales makes it feasible to evaluate $\zeta_{\rm conn}$
across the full BAO range for DESI-scale catalogues, where covariance
estimation requires $\mathcal{O}(10^3)$ mock realisations.

\paragraph{Parity violation}
\citet{hou23} reported a non-zero parity-odd 4PCF in BOSS DR12 using a
spherical-harmonic estimator.  The GRAMSCI configuration-space parity
estimator provides a complementary cross-check that is sensitive to the
full angular structure of the signal without assuming a harmonic
decomposition.  The deterministic treatment of degenerate tetrahedra
(Section~\ref{sec:parity}) is essential here: at the few-percent level the
odd channel is otherwise contaminated by floating-point noise whose sign
depends on the compiler and hardware.

\paragraph{BAO with higher-order statistics}
GRAMSCI supports anisotropic 3PCF and 4PCF (via $\mu$ binning for
redshift-space distortions), enabling joint 2PCF+3PCF+4PCF BAO fits that
extract additional signal-to-noise from the same catalogue.

%% ===================================================================
\section{Summary}
\label{sec:summary}

We have presented four additions to GRAMSCI:

\begin{enumerate}
\item \textbf{Merge-walk inner kernel.}  Replacing binary-search with
  two-pointer sorted-set intersection reduces the 3PCF inner loop from
  $O(m\log m)$ to $O(m)$ and the 4PCF inner loop --- which required two
  binary searches --- correspondingly further.  Measured speedups are
  1.5--2$\times$ for the 3PCF and 2.8--3.8$\times$ for the 4PCF at
  survey-relevant scales.  The key prerequisite --- adjacency lists sorted
  by node index --- was already guaranteed by the KD-tree construction and
  required no data-structure change.
\item \textbf{Parity-decomposed 4PCF.}  A direction-pixel byte stored per
  edge at graph-construction time allows the 4PCF to be decomposed into
  parity-even and parity-odd channels at negligible additional cost,
  enabling direct tests of parity violation.  Degenerate (zero-volume)
  tetrahedra are excluded from the odd channel deterministically.
\item \textbf{Python interface.}  A thin NumPy-based wrapper calls the
  Fortran engine via subprocess, handling all file I/O automatically and
  returning structured result objects.
\item \textbf{GPU port with out-of-core support.}  The complete query
  engine runs on a single GPU via OpenACC with measured speedups of
  $2.6\times$ (3PCF) to $9\times$ (4PCF) over a 64-thread node, and an
  automatic window-tiling scheme processes graphs far larger than device
  memory ($9\times10^9$ edges demonstrated on a 24\,GB card).  All kernels
  reproduce the CPU reference exactly.
\end{enumerate}

The code is publicly available at \url{https://github.com/csabiu/Gramsci}.

\vspace{5mm}

\software{GRAMSCI \citep{sabiu19}; \textsc{kdtree2} \citep{kennel04};
  NumPy \citep{harris20}; NVIDIA HPC SDK.}

%% ===================================================================
\appendix

\section{OpenACC implementation notes}
\label{app:openacc}

We record four \textsc{nvfortran}/OpenACC behaviours that silently destroy
performance or correctness in this class of kernel; each cost significant
debugging effort and none is prominent in the standard documentation.

\paragraph{Implicit gang--vector mapping}
A bare \texttt{!\$ACC PARALLEL LOOP GANG} over hubs lets the compiler map
the hub loop across gangs \emph{and} vector lanes, placing a different hub
on every lane of a warp.  Each lane then traverses a different adjacency
list: memory access is fully divergent and warp execution time is set by
the largest hub in the warp.  An explicit inner \texttt{!\$ACC LOOP
VECTOR} over the neighbour loop, with the hub loop on gangs only, changed
this kernel from $6\times$ slower than the CPU node to $2.6\times$ faster.

\paragraph{Private arrays are per-lane}
Arrays in a \texttt{PRIVATE} clause of a gang loop with
\texttt{VECTOR\_LENGTH} $> 1$ are allocated per vector lane, not per gang.
A 360\,KB per-hub scratch array at vector length 64 implies 23\,MB per
thread block --- far beyond the per-SM limit --- and fails at kernel launch
with \texttt{CUDA\_ERROR\_ILLEGAL\_ADDRESS}.  The workaround is a global
scratch array indexed by an explicit gang-slot loop.

\paragraph{Array reductions}
Reductions of multidimensional allocatable arrays do not reliably
propagate device results back to the host, and array reductions declared
on nested gang+vector loops re-reduce at the end of every inner-loop
instance.  Slot-strided partial arrays updated with \texttt{ATOMIC
UPDATE} and summed on the host are both correct and fast.

\paragraph{By-reference scalar arguments go stale}
Scalar arguments passed (by Fortran default) by reference into
\texttt{!\$ACC ROUTINE} device functions receive a one-shot device copy
that is \emph{not} refreshed on subsequent kernel launches: host-side
updates between launches are silently ignored.  In the out-of-core tiles
this manifested as out-of-bounds reads from the second tile onward, while
scalars used directly in kernel code (which become \texttt{firstprivate})
remained correct.  All scalar arguments of device routines must carry the
\texttt{VALUE} attribute.

\bibliography{gramsci_v2}
\bibliographystyle{aasjournalv7}

\end{document}
